\documentclass[11pt]{article}
\usepackage{amsmath,amssymb,amsthm}
\usepackage[margin=1in]{geometry}
\usepackage{hyperref}
\usepackage{tikz-cd}

\newtheorem{theorem}{Theorem}
\newtheorem{lemma}{Lemma}
\newtheorem{corollary}{Corollary}
\theoremstyle{definition}
\newtheorem{definition}{Definition}

\newcommand{\Z}{\mathbb{Z}}
\newcommand{\F}{\mathbb{F}}
\newcommand{\lean}[1]{\texttt{#1}}

\title{Unary maps are restrictions of linear maps and $k$-ary maps are restrictions of
homogeneous maps}
\author{Roman Bacik}
\date{November 22, 2026}

\begin{document}
\maketitle

\begin{abstract}
We study maps on a finite set $S$ under renaming elements of $S$ and extending $S$. We show that
every unary map is linear and every $k$-ary map is homogeneous with respect to such renamings and
embeddings.
\end{abstract}

\section{Introduction}

A cellular automaton updates all cells at once by a local rule $f\colon S^k\to S$ that reads the
states of $k$ neighbouring cells, and already two-neighbour automata can be
computation-universal~\cite{smith}. We look for the simplest functional description of an arbitrary rule $f$, while
permitting renaming and enlarging of the set of states: the states are embedded into a ring $R$
and $f$ becomes the restriction of a simple map on $R$. This changes neither the transitions
between the original states nor what the automaton computes.

This note continues~\cite{finlin} and~\cite{finbin}. In~\cite{finlin} it is shown that every map
on a finite set is represented by a linear map over some $\Z/m$; this is
Theorem~\ref{thm:unary} below. In~\cite{finbin} it is shown that this does not extend to binary
maps: for every $d$ there is a binary map on a finite set that is not represented by a
polynomial of degree $d$ over any commutative ring. Here we show that every map $S^k\to S$
is represented by a homogeneous polynomial over some prime field (Theorem~\ref{thm:binary}), and
that nothing simpler works in general (Theorem~\ref{thm:optimal}).

Let $S$ be a finite set and $R$ a commutative ring.

\begin{definition}
Let $k\ge 1$. An \emph{embedding} of a map $f\colon S^k\to S$ into a map $q\colon R^k\to R$ is an
injective map $j\colon S\to R$ with $j\circ f=q\circ j^k$, where
$j^k(x_1,\dots,x_k)=(j(x_1),\dots,j(x_k))$, that is, such that the square
\[
\begin{tikzcd}
S^k \arrow[r,"f"] \arrow[d,hook,"j^k"'] & S \arrow[d,hook,"j"] \\
R^k \arrow[r,"q"'] & R
\end{tikzcd}
\]
commutes. For $k=1$ the map $f\colon S\to S$ is \emph{unary}, and $j\circ f=q\circ j$. For $k=2$
the map $f$ is \emph{binary} and we write $f(x,y)$ and $q(u,v)$.
\end{definition}

If $f$ has an embedding into $q$, we say that $f$ is \emph{represented} by $q$ over $R$.

\begin{definition}
A map $q\colon R\to R$ is \emph{linear} if $q(u)=c\,u$ for some $c\in R$. A map
$q\colon R^2\to R$ is \emph{$r$-separable}, $r\ge 0$, if $q(u,v)=\sum_{k=1}^{r}g_k(u)\,h_k(v)$
for some maps $g_k,h_k\colon R\to R$. A polynomial in $k$ variables over $R$ is identified with
the map $R^k\to R$ it defines.
\end{definition}

If $j$ is an embedding of $f$ into $q$, then, identifying $S$ with $j(S)\subseteq R$, the map $q$ sends $j(S)^k$ into $j(S)$ and
agrees there with $f$, so $f$ is the \emph{restriction} of $q$ to $j(S)^k$ and $q$ extends $f$
to $R^k$. Thus Theorem~\ref{thm:unary} says that every unary map is a restriction of a linear
map, and Theorem~\ref{thm:binary} that every $k$-ary map is a restriction of a homogeneous
polynomial map.

\section{Results}

Throughout, $s=|S|$, $\Z/0=\Z$, and $[P]\in\{0,1\}$ is the truth value of a statement $P$; we write
$[P]_R$ when the arithmetic in $P$ is evaluated in a ring $R$.

\begin{theorem}\label{thm:unary}
Every map $f\colon S\to S$ is represented by a linear map over some $\Z/m$, $m\ge 1$.
\end{theorem}

\begin{theorem}\label{thm:binary}
Let $k\ge 1$. Every map $f\colon S^k\to S$ is represented by a homogeneous polynomial over some
$\F_p$, $p$ prime.
\end{theorem}

To show that Theorem~\ref{thm:binary} cannot be simplified we use the shift, equality and
inversion indicators. For $n\ge 2$ we identify $\Z/n$ with $\{0,\dots,n-1\}$ and let
\[
  S_n(x,y)=[x=y+1]_\Z,\qquad E_n(x,y)=[x=y],\qquad I_n(x,y)=[xy=1]_{\Z/n}\qquad(x,y\in\Z/n).
\]

\begin{theorem}\label{thm:optimal}
\begin{enumerate}
\item[(a)] Let $n\ge 2$ and $m\ge 0$. If $m=0$ or $m\ge n$, every binary map on an $n$-element
set is represented by an $n$-separable map over $\Z/m$. The maps $S_n$ and $E_n$ are not
represented by an $(n-2)$-separable map over $\Z/m$.
\item[(b)] For every prime $p$ and every commutative ring $R$, the map $I_p$ is not represented
over $R$ by a polynomial of total degree at most $p-1$.
\end{enumerate}
\end{theorem}

Part (b) is~\cite[Theorem~5]{finbin}; we include its proof for completeness. A polynomial of
total degree at most $d$ is $\sum_{k=0}^{d}u^kq_k(v)$, hence $(d+1)$-separable, so by (a) the
maps $S_n$ and $E_n$ are not represented by a polynomial of total degree at most $n-3$. Since the
term $q_0(v)$ does not depend on $u$, the proof of (a) gives one degree more:

\begin{corollary}\label{cor:degree}
Let $n\ge 2$ and $m\ge 0$. The maps $S_n$ and $E_n$ are not represented over $\Z/m$ by a
polynomial of total degree at most $n-2$.
\end{corollary}

\section{Proof of Theorem~\ref{thm:unary}}

A map $\varphi\colon S\to G$ into an abelian group is \emph{Sidon}~\cite{obryant} if
$\varphi(x)-\varphi(y)=\varphi(x')-\varphi(y')$ with $x\ne y$ and $x'\ne y'$ implies
$(x,y)=(x',y')$. A Sidon map is injective: if $\varphi(x)=\varphi(y)$ with $x\ne y$, then
$\varphi(x)-\varphi(y)=0=\varphi(y)-\varphi(x)$ while $(x,y)\ne(y,x)$.

\begin{lemma}\label{lem:pow2}
Let $i,i',l,l'\ge 0$ with $i\ne i'$ and $l\ne l'$. If $2^i-2^{i'}=2^l-2^{l'}$, then $i=l$ and
$i'=l'$.
\end{lemma}

\begin{proof}
Both sides have the same sign, so we may assume $i>i'$ and $l>l'$. Then
$2^{i'}(2^{i-i'}-1)=2^{l'}(2^{l-l'}-1)$ with odd second factors, so $i'=l'$, and then $i=l$.
\end{proof}

\begin{lemma}\label{lem:sidon}
For every $f\colon S\to S$ and $d_0\ge 0$ there are $n\ge 1$, $d\ge d_0$ and a Sidon map
$\varphi\colon S\to\Z/n$ with $\varphi(f(x))=d\,\varphi(x)$ for all $x\in S$.
\end{lemma}

This refines~\cite[Lemma~2.12]{finlin}, which uses the weights $i+1$ in place of $2^i$ and
obtains an injective map; the weights $2^i$ make it Sidon.

\begin{proof}
If $s\le 1$, take $n=1$, $d=d_0$ and $\varphi=0$. Let $s\ge 2$, identify $S=\{0,\dots,s-1\}$,
and let $A$ be the $s\times s$ matrix $A_{ik}=[f(i)=k]$, so that $(Az)_i=z_{f(i)}$. Put
$\beta_i=2^i$, $\chi(t)=\det(tI-A)\in\Z[t]$ and $w(t)=\operatorname{adj}(tI-A)\,\beta\in\Z[t]^s$.
From $(tI-A)\operatorname{adj}(tI-A)=\chi(t)I$ we get $Aw=tw-\chi(t)\beta$, that is,
\[
  w_{f(i)}(t)=t\,w_i(t)-\chi(t)\,\beta_i .
\]
The polynomial $\chi$ is monic of degree $s$. The diagonal entries of
$\operatorname{adj}(tI-A)$ are monic of degree $s-1$ and the other entries have degree at most
$s-2$~\cite[Lemma~2.7]{finlin},~\cite{horn}, so $w_i$ has degree at most $s-1$ and coefficient $\beta_i$ at $t^{s-1}$. For $i\ne i'$,
$l\ne l'$ and $(i,i')\ne(l,l')$, the polynomial $w_i-w_{i'}-w_l+w_{l'}$ has coefficient
$\beta_i-\beta_{i'}-\beta_l+\beta_{l'}\ne 0$ at $t^{s-1}$ by Lemma~\ref{lem:pow2}. Hence for
every large integer $t_0$ we have $0<|w_i(t_0)-w_{i'}(t_0)-w_l(t_0)+w_{l'}(t_0)|<\chi(t_0)$ for
all such quadruples. Fix such a $t_0\ge\max(d_0,2)$ and put $n=\chi(t_0)$, $d=t_0$ and
$\varphi(i)=w_i(t_0)\bmod n$. The displayed identity gives $\varphi(f(i))=d\,\varphi(i)$, and
$\varphi$ is Sidon because $n$ divides no nonzero integer of absolute value less than $n$.
\end{proof}

\begin{proof}[Proof of Theorem~\ref{thm:unary}]
Take $n$, $d$ and $\varphi$ from Lemma~\ref{lem:sidon} and put $j=\varphi$. Then $j$ is
injective and $j\circ f=q\circ j$ for $q(u)=d\,u$, so $j$ is an embedding of $f$ into the linear
map $q$ over $\Z/n$.
\end{proof}

\section{Proof of Theorem~\ref{thm:binary}}

\begin{lemma}\label{lem:power}
For every $f\colon S\to S$ and $d_0\ge 0$ there are a prime $p$, an integer $d\ge d_0$ and an
embedding $j\colon S\to\F_p^\times$ of $f$ into $q(u)=u^d$ over $\F_p$ such that the ratios
$j(x)/j(y)$, $x\ne y$, are pairwise distinct, that is, $j(x)j(y')=j(x')j(y)$ with $x\ne y$ and
$x'\ne y'$ implies $(x,y)=(x',y')$.
\end{lemma}

\begin{proof}
Take $n$, $d$ and $\varphi$ from Lemma~\ref{lem:sidon}, a prime $p\equiv 1\pmod n$~\cite{borevich},
and $h\in\F_p^\times$ of order $n$, which exists since $\F_p^\times$ is cyclic of order
$p-1$~\cite{lang}. Put $j(x)=h^{\varphi(x)}$. Then
$j(f(x))=h^{d\,\varphi(x)}=j(x)^d$. Since $k\mapsto h^k$ is injective on $\Z/n$, so is
$j$, and the ratios $j(x)/j(y)=h^{\varphi(x)-\varphi(y)}$, $x\ne y$, are pairwise
distinct because $\varphi$ is Sidon.
\end{proof}

\begin{proof}[Proof of Theorem~\ref{thm:binary}]
Apply Lemma~\ref{lem:power} to the diagonal $a\mapsto f(a,\dots,a)$ with $d_0=s^k$. This gives a
prime $p$, an integer $d\ge s^k$ and an embedding $j\colon S\to\F_p^\times$ with
$j(f(a,\dots,a))=j(a)^d$ and pairwise distinct ratios $j(a)/j(b)$, $a\ne b$. For
$x=(x_1,\dots,x_k)\in S^k$ put
\[
  P(x)=\Bigl(\frac{j(x_1)}{j(x_1)},\dots,\frac{j(x_k)}{j(x_1)}\Bigr)\in\F_p^k,\qquad
  V(x)=j(f(x))\,j(x_1)^{-d}.
\]
If $P(x)=P(x')$, then $V(x)=V(x')$. Indeed, $j(x_i)j(x'_1)=j(x'_i)j(x_1)$ for all $i$. If
$x_1=x'_1$, then $x=x'$. Otherwise $x_i\ne x'_i$ for every $i$, since $x_i=x'_i$ would give $j(x'_1)=j(x_1)$, and the distinct ratios give
$x_i=x_1$ and $x'_i=x'_1$, so $x$ and $x'$ are constant; then $V(x)=V(x')=1$ by the choice of $j$.

Let $T=P(S^k)$, so $|T|\le s^k\le d$, and let $V_a=V(x)$ for any $x$ with $P(x)=a\in T$. For
$a\ne b$ in $T$ choose an index $c=c(a,b)$ with $a_c\ne b_c$, and put
\[
  L_{a,b}(u)=\frac{u_c-b_c\,u_1}{a_c-b_c},\qquad B_a=\prod_{b\in T\setminus\{a\}}L_{a,b},\qquad
  q=\sum_{a\in T}V_a\,u_1^{\,d-|T|+1}B_a .
\]
Each $L_{a,b}$ is a linear form, so $q$ is homogeneous of degree $d$. At $u=j^k(x)$ we have
$L_{a,b}(u)=j(x_1)\,(P(x)_c-b_c)/(a_c-b_c)$, which is $j(x_1)$ if $P(x)=a$ and $0$ if $b=P(x)$.
Hence $B_a(j^k(x))=j(x_1)^{|T|-1}\,[P(x)=a]$ and
$q(j^k(x))=V(x)\,j(x_1)^d=j(f(x))$. So $j$ is an embedding of $f$ into $q$.
\end{proof}

\section{Proof of Theorem~\ref{thm:optimal}}

We use a rank criterion.

\begin{lemma}\label{lem:lift}
Let $m,r\ge 0$, let $P$ be an integer matrix, let $c\in\Z/m$ be nonzero, and suppose $cP=EH$
over $\Z/m$ for matrices $E$ and $H$ over $\Z/m$ with $r$ columns and $r$ rows respectively. Then
$P$ has rank at most $r$ over $\F_p$ for some prime $p$.
\end{lemma}

\begin{proof}
Lift $E$, $H$ and $c$ to integer matrices $E'$, $H'$ and an integer $c'$. Every entry of
$E'H'-c'P$ is divisible by $m$, and $m\nmid c'$, so there are a prime $p$ and $k\ge 1$ with
$p^k\mid m$ and $p^k\nmid c'$. Write $c'=p^ec''$ with $p\nmid c''$; then $e<k$, so
$p^{e+1}\mid m$ and $E'H'=p^eB$ with $B=c''P+pD$ for an integer matrix $D$. The
$(r+1)\times(r+1)$ minors of $E'H'$ vanish, since $E'H'$ factors through $\mathbb{Q}^r$; they are
$p^{e(r+1)}$ times the minors of $B$. So $B\equiv c''P$ has rank at most $r$ over $\F_p$, and
$c''$ is a unit modulo $p$.
\end{proof}

\begin{lemma}\label{lem:criterion}
Let $f\colon X\times Y\to Z$, $r\ge 0$, $x_0\in X$, finitely many $x_i\in X$ and $y_l\in Y$,
and $u\ne w$ in $Z$ with $f(x_0,y_l)=u$ and $f(x_i,y_l)\in\{u,w\}$ for all $i,l$, and suppose that
the matrix $P_{il}=[f(x_i,y_l)=w]$ has rank greater than $r$ over $\F_p$ for every
prime $p$. Then for no $m\ge 0$ are there an injective map $j\colon Z\to\Z/m$ and maps
$G_k\colon X\to\Z/m$ and $H_0,H_k\colon Y\to\Z/m$ with
$j(f(x,y))=H_0(y)+\sum_{k=1}^{r}G_k(x)H_k(y)$ for all $x,y$.
\end{lemma}

\begin{proof}
Put $c=j(u)-j(w)\ne 0$, $E_{ik}=G_k(x_0)-G_k(x_i)$ and $H'_{kl}=H_k(y_l)$. The terms
$H_0(y_l)$ cancel in $j(f(x_0,y_l))-j(f(x_i,y_l))$, so $(EH')_{il}=c\,P_{il}$. By
Lemma~\ref{lem:lift} the matrix $P$ has rank at most $r$ over some $\F_p$, a contradiction.
\end{proof}

\begin{proof}[Proof of (a)]
Let $f$ be a binary map on an $n$-element set $S$. Since $m=0$ or $m\ge n$, there is an
injective map $j\colon S\to\Z/m$. For $z\in S$ put
$g_z(u)=[u=j(z)]$, $h_z(j(y))=j(f(z,y))$ for $y\in S$, and $h_z(v)=0$ for
$v\notin j(S)$. Then $j$ is an embedding of $f$ into
$q(u,v)=\sum_{z\in S}g_z(u)\,h_z(v)$.

If $j$ is an embedding of $S_n$ into $\sum_{k=1}^{n-2}g_k(u)h_k(v)$, then
Lemma~\ref{lem:criterion} applies with $G_k(x)=g_k(j(x))$, $H_k(y)=h_k(j(y))$, $H_0=0$,
$r=n-2$, $x_0=0$, $x_i=i$ and $y_l=l-1$ for $1\le i,l\le n-1$, $u=0$ and $w=1$: then $P$ is the
$(n-1)\times(n-1)$ identity matrix, of rank $n-1>r$ over every $\F_p$, a contradiction. For $E_n$
use $y_l=l$ instead.
\end{proof}

\begin{proof}[Proof of Corollary~\ref{cor:degree}]
A polynomial of total degree at most $n-2$ has the form $\sum_{k=0}^{n-2}u^kq_k(v)$. If $j$ is an
embedding of $f=S_n$ or $f=E_n$ into it, then $j(f(x,y))=H_0(y)+\sum_{k=1}^{n-2}G_k(x)H_k(y)$ with
$G_k(x)=j(x)^k$, $H_k(y)=q_k(j(y))$ and $H_0(y)=q_0(j(y))$. This is excluded by
Lemma~\ref{lem:criterion} with the data of the proof of (a).
\end{proof}

For (b) let $\sigma_k(u_0,u_1)=\sum_{i=0}^{k-1}u_1^iu_0^{k-1-i}$, so that
$u_1^k-u_0^k=(u_1-u_0)\sigma_k(u_0,u_1)$, and let $q_{k_1k_2}$ be the coefficient of
$u^{k_1}v^{k_2}$ in $q$.

\begin{lemma}\label{lem:mixed}
If $q\in R[u,v]$ has total degree at most $N+1$, then for all $u_0,u_1,v_0,v_1\in R$
\[
  q(u_1,v_1)-q(u_0,v_1)-q(u_1,v_0)+q(u_0,v_0)=(u_1-u_0)(v_1-v_0)\sum_{k_1,k_2=1}^{N}
  q_{k_1k_2}\,\sigma_{k_1}(u_0,u_1)\,\sigma_{k_2}(v_0,v_1).
\]
\end{lemma}

\begin{proof}
The monomial $u^{k_1}v^{k_2}$ contributes $(u_1^{k_1}-u_0^{k_1})(v_1^{k_2}-v_0^{k_2})$, which
vanishes unless $k_1,k_2\ge 1$, and then $k_1,k_2\le N$.
\end{proof}

\begin{proof}[Proof of (b)]
Suppose $j$ is an embedding of $I_p$ into $q\in R[u,v]$ of total degree at most $p-1$. Put
$N=p-2$, $\delta=j(1)-j(0)$ and $U=(\Z/p)^\times$. For
$x,y\in U$, Lemma~\ref{lem:mixed} with $u_0=v_0=j(0)$, $u_1=j(x)$ and $v_1=j(y)$ gives, since
$I_p(0,y)=I_p(x,0)=I_p(0,0)=0$,
\begin{gather*}
  (j(x)-j(0))(j(y)-j(0))\,e(x,y)=j(I_p(x,y))-j(0)=\delta\,[xy=1]_{\Z/p},\\
  e(x,y)=\sum_{k_1,k_2=1}^{N}q_{k_1k_2}\,\sigma_{k_1}(j(0),j(x))\,\sigma_{k_2}(j(0),j(y)).
\end{gather*}
In matrix form, indexed by $U$, this reads $D\,\Phi Q\Phi^{\mathsf T}D=\delta F$, where
$D=\operatorname{diag}(j(x)-j(0))$, $\Phi_{x,k}=\sigma_k(j(0),j(x))$ for $1\le k\le N$,
$Q=(q_{k_1k_2})$, and $F$ is the permutation matrix of $x\mapsto x^{-1}$. Since
$\Phi Q\Phi^{\mathsf T}$ factors through $R^N$ and $N<p-1=|U|$, its determinant vanishes, while
$\det F=\pm1$. Taking determinants gives $\delta^{p-1}=0$. The entry $x=y=1$ gives
$\delta(1-\delta\,e(1,1))=0$; as $\delta\,e(1,1)$ is nilpotent, $1-\delta\,e(1,1)$ is a unit, so
$\delta=0$, contradicting injectivity.
\end{proof}

\section*{Acknowledgements}
All results are formalized in Lean~4~\cite{lean4} with Mathlib~\cite{mathlib}; the formalization
is available at~\cite{hombin}.
The formal verification was carried out with the help of Claude Opus~\cite{claude} and Aristotle~\cite{aristotle}.

\begin{thebibliography}{9}
\bibitem{finlin}
R.~Bacik.
\newblock A linear representation for functions on finite sets.
\newblock arXiv:2510.20167, 2025.

\bibitem{finbin}
R.~Bacik.
\newblock Embedding finite functions into low-degree polynomial functions over commutative
rings.
\newblock arXiv:2606.09045, 2026.

\bibitem{hombin}
R.~Bacik.
\newblock HomBin.
\newblock \url{https://github.com/roman3017/HomBin}, 2026.

\bibitem{smith}
A.~R.~Smith~III.
\newblock Simple computation-universal cellular spaces.
\newblock \emph{J.\ ACM} 18 (1971), 339--353.

\bibitem{obryant}
K.~O'Bryant.
\newblock A complete annotated bibliography of work related to Sidon sequences.
\newblock \emph{Electron.\ J.\ Combin.}, Dynamic Survey DS11, 2004.

\bibitem{horn}
R.~A.~Horn and C.~R.~Johnson.
\newblock \emph{Matrix Analysis}, 2nd edition.
\newblock Cambridge University Press, 2012.

\bibitem{borevich}
Z.~I.~Borevich and I.~R.~Shafarevich.
\newblock \emph{Number Theory}.
\newblock Academic Press, 1966.

\bibitem{lang}
S.~Lang.
\newblock \emph{Algebra}, revised 3rd edition.
\newblock Springer, 2002.

\bibitem{mathlib}
The mathlib Community.
\newblock The Lean mathematical library.
\newblock In \emph{Proceedings of the 9th ACM SIGPLAN International Conference on Certified
Programs and Proofs (CPP 2020)}, pages 367--381, 2020.

\bibitem{lean4}
L.~de~Moura and S.~Ullrich.
\newblock The Lean 4 theorem prover and programming language.
\newblock In \emph{Automated Deduction -- CADE 28}, Lecture Notes in Computer Science 12699,
pages 625--635, 2021.

\bibitem{claude}
Anthropic.
\newblock Claude Opus [Large language model], 2026.
\newblock \url{https://claude.ai}.

\bibitem{aristotle}
The Harmonic Team.
\newblock Aristotle: IMO-level automated theorem proving.
\newblock arXiv:2510.01346, 2025. \url{https://aristotle.harmonic.fun}.
\end{thebibliography}

\end{document}
